% ── ABSTRACT (~300 words) ─────────────────────────────────────────────────────

Three independent intellectual traditions --- modern quantum physics, the
process philosophy of Alfred North Whitehead, and the Bahá'í writings of
Bahá'u'lláh and 'Abdu'l-Bahá --- converge on the same structural
description of reality: that existence is constituted by relationship, that
properties emerge only through relational interaction, and that the force
sustaining those relationships is, in each tradition's vocabulary, love.

This paper makes three related arguments.
First, substance ontology is \emph{formally incompatible} with quantum mechanics:
Bell's theorem \citep{Bell1964}, combined with the experimental violation of
its inequalities \citep{Aspect1982,Hensen2015}, makes local realism untenable
and undermines the substance-ontological assumption that systems carry
intrinsic properties prior to measurement, making Whitehead's process ontology
not merely a permitted interpretation but arguably the most systematically
developed internally consistent metaphysical framework available for quantum
reality.
Second, the Bahá'í writings constitute an \emph{independent prior instantiation}
of process ontology, articulating its essential claims decades before
\citet{Whitehead1929} named the framework philosophically and roughly a century
before physics confirmed it experimentally.
Third, these two findings together establish a \emph{triadic convergence}:
three traditions built by three different civilizations, using three different
methodologies, across 170 years, detect the same feature of reality.

The convergence is formalized through six structural isomorphisms (relational
ontology, nonlocal interconnection, observer-dependence, decoherence, holographic
order, and motion as life) and the Relational Coherence Theorem (\RCT),
developed in a companion physics paper~\citep{Adams2026Physics}, which translates
the philosophical claim into testable quantum-mechanical predictions.
The argument employs an overlapping consensus architecture: each tradition
validates its own thesis independently; the convergence is not a merged view
but an independent confirmation by separate instruments of the same underlying
structure of reality. In each tradition's vocabulary the relational principle
sustaining this structure is named \emph{love} --- \mahabbat in the Bahá'í
lexicon, the creative advance in Whitehead, and the nonlocal coherence
revealed by entanglement in quantum theory --- and the convergence on this
single substantive identification is itself the central finding.
